\begin{problem}[第30届全国中学生物理竞赛决赛][超导圆环在磁棒上方运动]{13}
如图，处于超导态、半径为 $r_0$ 、质量为 $m$ 、自感为 $L$ 的超导细圆环处在竖直放置的柱形磁棒上方，其对称轴与磁棒的对称轴重合。圆环附近的磁场具有柱对称性，磁感应强度 $\pmb{B}$ 可用一个竖直分量 $B_z = B_0(1 - 2\alpha z)$ 和一个径向分量 $B_r = B_0\alpha r$ 近似地描述。这里 $B_0$ 、 $\alpha$ 为大于零的常量， $z$ 、 $r$ 分别为竖直和径向位置坐标。在 $t = 0$ 时刻，环心的坐标为 $z = 0$ 、 $r = 0$ ，环上电流为 $I_0$ （规定圆环中电流的正向与 $z$ 轴正向满足右手规则）。此时把圆环释放，圆环开始向下运动，其对称轴仍然保持竖直。处于超导态的超导细圆环具有这样的性质：穿过超导细圆环的磁通量保持不变。
\begin{subproblem}
圆环作何种运动？给出环心的 $z$ 坐标与时间的依赖关系。
\end{subproblem}
\begin{subproblem}
求 $t$ 时刻圆环中电流 $I$ 的表达式。
\end{subproblem}
\end{problem}